\subsection{A family with the path structure}\label{sec:mechanism}

The benchmark models do not tell whether the cuts can supply strength that
SCIP lacks where the theory predicts a gap. We therefore built a family with
the structure of Section~\ref{sec:interleave}. An instance with $n$ copies has
variables $x_i,y_i,z_i\in[0,1]$ and $t_i$, the rows
$\Phi^{(i)}(x_i,y_i,z_i)-t_i\le0$, where $\Phi^{(i)}$ is the
function~\eqref{eq:family} with $\omega_A=\omega_C=1$, one coupling row
$\sum_iy_i\le0.8n$, and the objective $\min\sum_it_i$. For each copy, four
distinct values are drawn from $\{k/64:k=0,\dots,48\}$ and split into two
interleaving two-point sets $A_i$ and $C_i$ with random orientation. All
coefficients are dyadic and exactly representable. Every optimal $y_i$ is a
midpoint of two values of at most $3/4$, so the coupling row does not bind, and
by Theorem~\ref{thm:interleave} the optimum is $\sum_i\delta_i^2/2$. Glued exact
pair hulls give the bound $0$ for every copy, so the part of the gap that only
the joint block closes is the whole optimum. Campaign~3 used five instances for
each $n\in\{10,20,40,80\}$ (seeds 0--4); campaign~4 adds five fresh instances
per $n$ (seeds 5--9), which no earlier run had used. The randomized family was
adopted after a pilot with native SCIP alone had shown that the deterministic
family, with every copy equal to the example of
Proposition~\ref{prop:path-gap}, is easy for SCIP: it solved $160$ copies in
19\,s.

Because the objective is a sum of block functions and the coupling row does not
bind, the $n$ cuts $t_i\ge\min\Phi^{(i)}$, which are the exact supports of the
blocks in the directions of their own rows, make the root bound equal to the
optimum. This is the most favorable case for block cuts. Part~C4 removes it: it
replaces the coupling row of each fresh instance by $\sum_iy_i\le c$, where $c$
is half the sum of the smallest unconstrained minimizers $y_i^\ast$, rounded down
to a multiple of $1/64$. The row then binds, the block minima no longer give the
optimum, and the cuts must be found in other directions. For C4 we know two
reference values exactly: the optimum, from a convex mixed-integer quadratic
reformulation solved by Gurobi and certified by an exact branch and bound; and
the best root bound that cuts on the blocks $(x_i,y_i,z_i)$ can give, which by
Theorem~\ref{thm:closure} and Corollary~\ref{cor:one-row} is
$\min\{\sum_i\operatorname{vex}\varphi_i(y_i):\sum_iy_i\le c,\ 0\le y\le1\}$
with $\varphi_i(y)=\dist(y,A_i)^2+\dist(y,C_i)^2$, computed with the exact convex
envelopes of the piecewise quadratic functions $\varphi_i$.
Table~\ref{tab:mechanism} summarizes the results; Tables~\ref{tab:path-detail}
and~\ref{tab:path-fresh} in Appendix~\ref{app:instances} give them per
instance.

\begin{table}[t]
\centering\footnotesize
\caption{The path family: median fraction of SCIP's default root gap closed,
$(z_{\text{mode}}-z_{\text{base}})/(z^\ast-z_{\text{base}})$, for each $n$ (five
instances each), and instances solved within 300\,s (of 20). ``Remainder'' and
``whole row'' are the first directions of the separator
(Section~\ref{sec:blocks}); ``mech.'' and ``wide'' are the limits of
Table~\ref{tab:modes}. The row ``glued pair hulls'' is the analytic bound~0 of
Theorem~\ref{thm:interleave}, which no run computes.}\label{tab:mechanism}
\begin{tabular}{@{}llrrrrr@{}}
\toprule
& & \multicolumn{4}{c}{Root gap closed (median)} & \\
\cmidrule(lr){3-6}
Instances & Mode & $n=10$ & $n=20$ & $n=40$ & $n=80$ & Solved\\
\midrule
seeds 0--4 & SCIP default & --- & --- & --- & --- & 9\\
 & SCIP, \code{checkvarlocks} off & 0.95 & 0.89 & 0.93 & 0.91 & 10\\
 & SCIP, disabled separators on & 0.64 & 0.58 & 0.66 & 0.59 & 10\\
 & Gurobi & & & & & 7\\
 & glued pair hulls & 0.95 & 0.89 & 0.94 & 0.92 & \\
 & remainder, mech.\ (campaign 3) & 0.35 & 0.25 & 0.19 & 0.22 & 9\\
 & remainder, wide & 0.86 & 0.87 & 0.89 & 0.90 & 13\\
 & whole row, mech.\ (post hoc) & 0.44 & 0.44 & 0.43 & 0.46 & 10\\
 & whole row, wide (post hoc) & 1.00 & 1.00 & 1.00 & 1.00 & 20\\
\midrule
seeds 5--9 & SCIP default & --- & --- & --- & --- & 9\\
 & SCIP, disabled separators on & 0.57 & 0.64 & 0.61 & 0.66 & 9\\
 & Gurobi & & & & & 5\\
 & glued pair hulls & 0.93 & 0.92 & 0.93 & 0.92 & \\
 & remainder, mech. & 0.32 & 0.31 & 0.24 & 0.21 & 10\\
 & whole row, wide & 1.00 & 1.00 & 1.00 & 1.00 & 20\\
\midrule
binding row & SCIP default & --- & --- & --- & --- & 10\\
(seeds 5--9) & SCIP, disabled separators on & 0.29 & 0.60 & 0.68 & 0.65 & 13\\
 & Gurobi & & & & & 6\\
 & remainder, wide & 0.98 & 0.96 & 0.95 & 0.93 & 20\\
 & whole row, wide & 0.997 & 0.99 & 0.98 & 0.98 & 20\\
\bottomrule
\end{tabular}
\end{table}

\paragraph{SCIP's root bound and the pair-hull level.} SCIP's default root
bound is weak: on the instances of campaign~3 it lies between $-0.035$ and
$-0.013$ per copy, while each copy contributes at least $1/8192$ to the optimum. The
cause is a presolve step: $\Phi^{(i)}$ is concave in $x_i$ and in $z_i$, so
SCIP's implicit-discreteness detection \citep[Section~4.2.7]{SCIP8report} makes
these variables binary, and SCIP then closes the gap by branching. With this
step disabled (\code{constraints/nonlinear/checkvarlocks = d}), SCIP's root bound
rises to just below~0, the bound of glued exact pair hulls, on every instance
of campaign~3 (Table~\ref{tab:mechanism}); SCIP's semidefinite minor cuts were
productive in all of these runs and in none of the default runs of campaign~4. Native SCIP thus reaches the
pair-hull level when presolve does not interfere, and no SCIP setting closed
any part of the gap between this level and the optimum, which is the gap that
Section~\ref{sec:composition} attributes to joint blocks. SCIP's disabled
separators, of which only the intersection cuts were productive, closed about
60\% of the default root gap. Gurobi~13 solved the instances with $n=10$ and
two with $n=20$; on $n=40$ and $n=80$ its gap after 300\,s was 53 to 121\%.

\paragraph{The prospective separator and the attribution.} The separator of
campaign~3 improved the root bound on all 20 instances but closed only a median
of 19--35\% of the root gap, less than the pair-hull level, and it solved the
same 9 instances as SCIP; on these, it reduced the number of nodes and the
shifted geometric mean time from 10.3\,s to 6.7\,s. After seeing these results
we found two causes. The first directions of the separator bound only the
nonlinear remainder of a row (Section~\ref{sec:blocks}), whereas the exact
support of the whole row is the block minimum, and the $n$ cuts
$t_i\ge\min\Phi^{(i)}$ give the optimum. In addition, the first blocks of a
callback take up to four cuts each, so with $n$ cuts per callback the later
blocks receive none. A post hoc diagnostic, designed after these results and
reported separately, tried the whole-row direction first, with the same limits
and with four times as many cuts per callback and in total and twice as many
support calls. Before it, we had
checked the wide limits by hand on two of the 20 instances. The whole-row
direction alone raised the median closure to 43--46\% and solved one more
instance; together with the wide
limits it closed the root gap on all 20 instances, and all 20 were solved, 12
of them at the root node. Campaign~4 added the missing cell of this $2\times2$
design prospectively: with the remainder directions and the wide limits, the
separator closed 86--90\% of the gap, close to the pair-hull level and above
it on 7 of the 20 instances, and solved 13. The campaign-4 runs had a less loaded
host than those of campaign~3, but this does not explain the difference: every
run of the four cells stopped at its cut cap, so the root bounds do not depend
on time, and the four additional instances were solved in 9 to 168\,s. The
limits therefore account for most of the closure, and the whole-row
directions for the part beyond the pair-hull level.

\paragraph{Prospective test on fresh instances.} On the 20 instances with
seeds 5--9, the whole-row variant with wide limits, fixed before these
instances were generated, closed the root gap on every instance: its root
bound was within $2\cdot10^{-7}$ of the optimum. It solved all 20 in 3 to 33\,s,
with a median of one node and at most 38, against 9 for SCIP, 9
for SCIP with its disabled separators, 10 for the separator with remainder directions and mechanism limits, and 5 for
Gurobi. On the nine instances that SCIP solved, the total time was comparable
(median ratio 0.86), while SCIP's own time fell to a median of 4\% of its time
in the baseline; over all 20 runs, 96\% of the time was spent in the Python
separator. The separator with remainder directions and mechanism limits
closed 21--32\% of the root gap, as in campaign~3.

\paragraph{A binding coupling row (Part C4).} When the coupling row binds,
the block minima no longer give the optimum, and the cuts must have the right
slopes in $y_i$. Both separator configurations with wide limits solved all 20
instances, in 3 to 134\,s, against 10 for SCIP, 13 for SCIP with its disabled
separators and 6 for Gurobi; only the cut modes solved any instance with
$n=80$. On the ten instances that SCIP solved, all with $n\le20$, the cut
modes were slower by median factors of 3.2 and 3.5, because of the
separator's time. At the root they closed a median of 95\% (remainder directions) and 99\%
(whole-row directions) of SCIP's root gap. Here the reference that matters is
the best root bound that cuts on the blocks can give, which equals the optimum
on 13 of the 20 instances and lies at most $6.5\cdot10^{-4}$ below it on the
others; no root run reached it, and the remaining distance grew with $n$, to
between 0.016 and 0.029 for $n=80$ with whole-row directions. Every cut-mode run
stopped at its cap of $16n$ cuts, so a larger cap would probably close more. The
cuts that did the work came from the LP direction search: 99.9\% (remainder directions) and
94\% (whole row) of the cuts were LP directions, and almost all of them have a
nonzero coefficient of $y_i$, mostly with the negative slope that the binding
row requires; the whole-row cuts alone give only $\sum_i\min\varphi_i$, which is
below SCIP's root bound on half of the instances. The direction search can
therefore find useful non-row directions, but on these instances it needed
about $16$ cuts per block to come close to the block closure, and none of the
cuts is an envelope cut in $(y_i,t_i)$ alone.

These families were constructed to have the structure of
Section~\ref{sec:composition}, and their objectives are sums of block
functions, which favors block cuts. They show that exact block support
supplies strength that neither SCIP, with or without its disabled separators,
nor Gurobi obtains at the root, that this strength reaches the solver only if
the direction search tries the right directions and enough of them, and that
on such instances it can turn time-outs into solves at or near the root. They
are not evidence about other models.
